\section{Introduction}
\label{sec:introduction}

\if{
Quantum key distribution (QKD) promises information-theoretic security grounded in the laws of quantum mechanics \cite{Bennett1984, Ekert1991}. While canonical protocols such as BB84 \cite{Bennett1984, Bennett1984Original, Bennett1992, Shor2000, Mayers2001, Scarani2009, Bennett1992_exp} rely on transmitting and measuring single-qubit states, recent theory and experiments have opened alternative cryptographic primitives. Among these is quantum energy teleportation (QET)—the remote extraction of energy using only local operations and classical communication (LOCC) on a shared entangled ground state \cite{Hotta2011}. QET does not teleport a quantum state but rather the expectation value of a local observable. Demonstrations on superconducting hardware \cite{Ikeda2023} have established feasibility and, more broadly, motivated extensions from energy to other globally conserved quantities such as charge and current \cite{IkedaTeleportingCharge}—an approach we refer to as quantum observable teleportation.

Building on this idea, it was proposed to encode key bits in the sign of a teleported observable’s expectation value, controlled by a single classical bit chosen by Alice
\cite{QKDbyQET}. Such a mechanism offers a general framework for secure signaling: Alice deterministically controls a binary outcome at Bob’s location while using only LOCC. Because the required resource is a carefully prepared many-body ground state, the approach is naturally suited to applications where unconditional security is prioritized over throughput—precisely the trade-off emphasized in QKD. 

In this work we develop and analyze a QKD protocol based on charge teleportation. We study two implementations using transverse-field Ising models (TFIM): a star-coupled architecture in which Alice interacts with all other spins, and a nearest-neighbor one-dimensional chain with finite separation between Alice and Bob. Using numerical simulations and Qiskit-level circuit models, we compare charge vs energy teleportation and find that charge provides superior bit-symmetry, statistical stability, and scalability with system size, making it a strong candidate for practical QKD. We also quantify robustness to classical and quantum noise channels, identifying operating regimes in which secure key distribution remains viable.

\paragraph{Organization:} Section~\ref{sec:protocol} reviews the minimal observable-teleportation for a generic system and reintroduces the protocol designed in QKD by QET~\cite{QKDbyQET}. Section~\ref{sec:hamiltonians} presents the analysis for the two Hamiltonians we use for our model and the optimal control angles. Section~\ref{sec:numerical_sim} discusses numerical simulations of the protocol for both models, while Section~\ref{sec:qiskit_sim} provides validation of the numerical simulations using quantum circuits and presents results from execution on real quantum hardware. Section~\ref{sec:noise-errors} analyzes the protocol's resilience to different noise and error models, followed by a discussion and outlook in Section~\ref{sec:summary}.

}\fi

Quantum Key Distribution (QKD) promises unconditionally secure communication by leveraging the fundamental principles of quantum mechanics \cite{Bennett1984, Ekert1991}. While traditional protocols like the BB84 protocol~\cite{Bennett1984, Bennett1984Original, Bennett1992, Shor2000, Mayers2001, Scarani2009, Bennett1992_exp, Pirandola2020} rely on the transmission and measurement of single-qubit states, recent theoretical and experimental advances have opened new avenues for cryptographic tasks. Among these is Quantum Energy Teleportation (QET), a protocol that enables the extraction of energy from a distant location by performing only local operations and classical communication (LOCC) on a shared entangled ground state, without violating local energy conservation \cite{Hotta2011, Amico2008}. This process does not teleport a quantum state, but rather the expectation value of a local observable.

Recent work has demonstrated the experimental feasibility of QET on superconducting quantum hardware \cite{Ikeda2023} and generalized the protocol beyond energy to other globally conserved quantities, such as charge and current \cite{IkedaTeleportingCharge}. This generalization, termed Quantum Observable Teleportation, provides a versatile toolkit for manipulating local properties of many-body systems. Building on this, it was recently proposed that such a mechanism could form the basis of a novel QKD protocol, where information is encoded in the sign of the teleported observable's expectation value, controlled by a classical bit choice made by the sender (Alice) \cite{QKDbyQET}. More broadly, this mechanism establishes a general framework for secure quantum signaling, where Alice can deterministically control a binary outcome at Bob's location. However, given that the protocol's resource is a meticulously prepared many-body ground state, its practical application is best suited for tasks where unconditional security is prioritized over high data rates. QKD is the canonical example of such a task and thus forms the central focus of our investigation.

Focusing on this primary application, we develop and analyze a QKD protocol based on the teleportation of a local charge observable. We investigate the protocol's performance under two distinct interaction models for the underlying many-body system: a star-coupled Transverse Field Ising Model (TFIM), $H^{(1)} = J\sum_{k=1}^{N}X_{0}X_{k}+h\sum_{k=0}^{N}Z_{k}$, where a central qubit (Alice) interacts with all others, and a 1D nearest-neighbor TFIM, $H^{(2)} = J\sum_{k=1}^{N}X_{k-1}X_{k}+h\sum_{k=0}^{N}Z_{k}$. Through numerical simulations and Qiskit-based quantum circuit analysis~\cite{Amir_QKD_by_Charge_Teleportation}, we compare the efficacy and robustness of charge teleportation against energy teleportation for these models. Our findings show that charge teleportation exhibits superior symmetry, stability against statistical noise, and scalability with system size, making it a more promising candidate for practical QKD implementations. Furthermore, we analyze the protocol's resilience to various noise channels, identifying the operational regimes where secure key distribution remains viable.

\paragraph{Organization:} Section~\ref{sec:protocol} reviews the minimal observable-teleportation for a generic system and reintroduces the protocol designed in~\cite{QKDbyQET}. Section~\ref{sec:hamiltonians} presents the analysis for the two Hamiltonians we use for our model. Section~\ref{sec:numerical_sim} discusses numerical simulations of the protocol, while Section~\ref{sec:qiskit_sim} provides validation of the numerical simulations using quantum circuits. Section~\ref{sec:noise-errors} analyzes the protocol's resilience to different noise and error models. Section~\ref{sec:security_analysis} then provides a cryptographic security analysis \cite{Renner2005}, quantifying the secret key rate under noise. The paper concludes with a summary and outlook in Section~\ref{sec:summary}.
